\section{平移差}
\label{sec:besov-translation-differences}

导数描述无穷小尺度上的变化率。若函数在某些点没有可积的导数，我们仍可以退一步，比较函数与它自身平移一小步$h$以后的差别，并观察这个差别怎样随$h$缩小。本节先用一个周期阶跃说明，平移差能读出导数条件看不到的分数次规律；再区分同一尺度上变化的大小与位置，以及不同尺度之间的汇总；最后用高阶差分和连续模把这些想法变成可以计算的量。

本章记$\Delta_hf(x)=f(x+h)-f(x)$，$f$为周期$1$的函数，$L^p$范数取在一个周期上。平移不改变$L^p(\mathbb T)$范数，这一点在后面反复使用。

\subsection{阶跃函数的平移差}
\label{sec:besov-step-differences}

取单位周期阶跃$f$：在$[0,1/2)$上为$1$，在$[1/2,1)$上为$0$，再按周期$1$延拓。它在$x=1/2$与$x=0$处跳跃，由第~\ref{chap:sobolev-spaces}章的论证，它在包含跳跃点的区间上没有弱导数，不满足一阶弱导数可积的条件。

我们直接计算平移差。设$0<h\leq1/2$。$\Delta_hf(x)\neq0$当且仅当$x$与$x+h$恰有一个落在$[0,1/2)$的周期复制中：$x\in[1/2-h,1/2)$时$f(x)=1$、$f(x+h)=0$；$x\in[1-h,1)$时$f(x)=0$、$f(x+h)=1$。在其余位置两者相同。所以$|\Delta_hf|$在总长度为$2h$的集合上等于$1$，其余为零，于是
\begin{equation}
\lVert\Delta_hf\rVert_p=(2h)^{1/p},\qquad 0<h\leq\frac12,\ 1\leq p<\infty .
\label{eq:besov-step-difference}
\end{equation}
$h<0$时把$x$换成$x+h$，得到同一结果，只需把$h$换成$|h|$。特别地，$\lVert\Delta_hf\rVert_2=(2|h|)^{1/2}$。

把它与光滑函数比较。对$g(x)=\sin(2\pi x)$，和差化积给出$\Delta_hg(x)=2\sin(\pi h)\cos(2\pi x+\pi h)$，所以$\lVert\Delta_hg\rVert_2=\sqrt2\,|\sin(\pi h)|\leq\sqrt2\,\pi|h|$。一般地，若周期函数$g$连续可微，由$\Delta_hg(x)=\int_0^hg'(x+t)\,\mathrm dt$与积分的Minkowski不等式，$\lVert\Delta_hg\rVert_p\leq|h|\,\lVert g'\rVert_p$。光滑函数的平移差与步长成正比，阶跃函数的平移差在$L^2$中只与步长的$1/2$次幂成正比。

这说明在“有可积导数”与“毫无规律”之间，还有可以检查的尺度规律：变化量是步长的分数次幂。阶跃函数虽然没有弱导数，它在$L^2$中仍有$1/2$阶的规律性。这种按步长幂次度量的规律，正是本章要刻画的光滑性。

\subsection{位置分布与尺度汇总}
\label{sec:besov-distribution-aggregation}

只给出一个幂次还不够，因为两个函数即使在小尺度上的变化有同样的上界，变化集中在哪些位置、又在多少个尺度上持续出现，仍可能不同。我们分别看这两个方面。

\paragraph{同一尺度内的位置分布}
式\eqref{eq:besov-step-difference}显示，同一个阶跃函数在不同积分指数下表现出不同的幂次：$p=1$时为$1$，$p=2$时为$1/2$，$p\to\infty$时趋于$0$。原因是变化集中在长度为$2h$的小集合上，幅度却保持为$1$。$p$越大，积分越看重少数大值，集中的变化就显得越剧烈。光滑的$\sin(2\pi x)$则不同，它的变化分散在整个周期上，对每个$p$都是$1$次。所以积分指数$p$记录了同一尺度上变化在位置上的分布。前章的积分指数已经起过类似作用，例如$W^{1,p}$是否包含某个幂函数取决于$p$；本章把这一作用与后面的尺度汇总分开规定。

\paragraph{不同尺度之间的汇总}
为看清尺度汇总，考虑单一频率的分量$u_j(x)=2^{-js}\cos(2\pi2^jx)$，$0<s<1$。与正弦函数的计算相同，
\[
\lVert\Delta_hu_j\rVert_2=2^{-js}\sqrt2\,|\sin(\pi2^jh)| .
\]
当$h$约为$2^{-j}$时，它与$2^{-js}$同阶；当$h\ll2^{-j}$时，它约为$\sqrt2\pi2^{j(1-s)}h$，小得多。所以第$j$个分量主要在尺度$h\approx2^{-j}$处“活跃”。

把这些分量按权重叠加：对$0\leq b_j\leq1$，令$f_b=\sum_{j\geq1}b_ju_j$。各分量频率不同，在$L^2$中正交，所以$\lVert\Delta_hf_b\rVert_2^2=\sum_jb_j^2\,2^{-2js}\cdot2\sin^2(\pi2^jh)$。用$|\sin\theta|\leq\min(1,|\theta|)$，把$2^jh\leq1$与$2^jh>1$的项分开求几何和，可得对每个这样的权重序列都有$\lVert\Delta_hf_b\rVert_2\leq C_sh^s$，$C_s$只依赖$s$。另一方面，取$h=2^{-J-1}$，单独第$J$项就给出$\lVert\Delta_hf_b\rVert_2\geq\sqrt2\,b_J2^{-Js}=\sqrt2\,b_J(2h)^s$。所以在尺度$2^{-J}$附近，$h^{-s}\lVert\Delta_hf_b\rVert_2$的大小主要由$b_J$及相邻几个$b_j$决定。

比较$b_j=1$与$b_j=1/j$两种权重。两个函数满足同样的统一上界$\lVert\Delta_hf_b\rVert_2\leq C_sh^s$：前者在每个二进尺度上都达到这个量级，后者在越来越细的尺度上逐渐减弱。统一上界只要求各尺度的贡献$(b_j)$有界，看不出这一区别。若进一步问这些贡献作为一个序列有多大，例如是否平方可和，两者就分开了：$(1/j)$平方可和，$(1)$不可和。所以除了同一尺度内的度量方式，还需单独规定不同尺度上的变化怎样汇总，第~\ref{sec:besov-norms}节的第二个指数承担这项工作。

\subsection{有限差分}
\label{sec:besov-finite-differences}

一阶平移差只能读出不超过一阶的规律：光滑函数的一阶差分也只与$h$成正比，读不出更高的光滑性。为读取更高阶的变化，把平移差迭代起来。对正整数$r$，定义$r$阶\term{有限差分}（finite difference）$\Delta_h^r=\Delta_h(\Delta_h^{r-1})$，$\Delta_h^1=\Delta_h$。展开得
\[
\Delta_h^rf(x)=\sum_{i=0}^r(-1)^{r-i}\binom ri f(x+ih).
\]
例如$\Delta_h^2f(x)=f(x+2h)-2f(x+h)+f(x)$。

高阶差分的作用在于抵消低阶多项式。$\Delta_h$把次数为$m$的多项式变成次数不超过$m-1$的多项式，所以$\Delta_h^r$消去所有次数小于$r$的多项式：$\Delta_hx=h$，$\Delta_h^2x=0$，$\Delta_h^2x^2=2h^2$。周期函数在整体上不是多项式，但若$f$在区间$[x,x+rh]$上与一个次数小于$r$的多项式重合，则$\Delta_h^rf(x)=0$。因此$r$阶差分忽略函数中局部的低阶多项式部分，只读出超出它的变化。

对光滑函数可以写出精确的积分表示。若周期函数$f$有$r$阶连续导数，$h>0$，对$r$归纳可得
\[
\Delta_h^rf(x)=\int_{[0,h]^r}f^{(r)}(x+t_1+\cdots+t_r)\,\mathrm dt_1\cdots\mathrm dt_r .
\]
由积分的Minkowski不等式，$\lVert\Delta_h^rf\rVert_p\leq h^r\lVert f^{(r)}\rVert_p$。另一方面，每作一次差分范数至多加倍，所以对任意$f\in L^p(\mathbb T)$都有$\lVert\Delta_h^rf\rVert_p\leq2^r\lVert f\rVert_p$。

\begin{example}[折线的二阶差分]
\label{ex:besov-triangle-wave}
令$g$为周期$1$的三角波：在$[0,1)$上$g(x)=|x-1/2|$。它连续，斜率处处为$\pm1$，在$x=1/2$和整数点有折点。一阶差分$|\Delta_hg|$在远离折点处等于$h$，所以$\lVert\Delta_hg\rVert_p$与$h$同阶，读不出比一阶更多的东西。二阶差分则只在$[x,x+2h]$包含折点时非零，这样的$x$组成总长度为$4h$的集合；在那里，由折线的凸性或凹性，$|\Delta_h^2g(x)|\leq2h$。所以$\lVert\Delta_h^2g\rVert_p\leq2h(4h)^{1/p}$，是$1+1/p$阶的。二阶差分看见了一阶差分看不见的规律：三角波除去折点处的集中变化以外，处处是线性的。
\end{example}

\subsection{连续模}
\label{sec:besov-modulus}

差分依赖一个具体的步长$h$。为了得到只依赖尺度的量，对所有不超过给定尺度的步长取最坏情形。

\begin{definition}[连续模]
\label{def:besov-modulus}
设$1\leq p\leq\infty$，$r$为正整数，$f\in L^p(\mathbb T)$。对$t>0$，称
\[
\omega_r(f,t)_p=\sup_{0<|h|\leq t}\lVert\Delta_h^rf\rVert_p
\]
为$f$在$L^p$中的$r$阶\term{连续模}（modulus of continuity）。
\end{definition}

它量化了步长不超过$t$时的统一变化。由定义，$\omega_r(f,t)_p$随$t$单调不减；由上一小节，$\omega_r(f,t)_p\leq2^r\lVert f\rVert_p$，对$r$阶连续可微的$f$还有$\omega_r(f,t)_p\leq t^r\lVert f^{(r)}\rVert_p$。差分对$f$线性，再由Minkowski不等式，$\omega_r(f+g,t)_p\leq\omega_r(f,t)_p+\omega_r(g,t)_p$。对阶跃函数，式\eqref{eq:besov-step-difference}给出$\omega_1(f,t)_p=(2t)^{1/p}$，$0<t\leq1/2$。

为什么要用高于待刻画阶数的差分？连续模的增长不能比$t^r$更慢，除非函数是常数。在$L^2$中这可以直接证明。记$\widehat f(k)=\int_0^1f(x)e^{-2\pi\mathrm ikx}\,\mathrm dx$。把第~\ref{chap:fourier-transform}章周期$2\pi$的Parseval等式换元到周期$1$，得$\lVert f\rVert_2^2=\sum_k|\widehat f(k)|^2$；平移对应于系数乘以$e^{2\pi\mathrm ikh}$，而$|e^{2\pi\mathrm ikh}-1|=2|\sin(\pi kh)|$，所以
\begin{equation}
\lVert\Delta_h^rf\rVert_2^2=\sum_{k\in\mathbb Z}\bigl(2|\sin(\pi kh)|\bigr)^{2r}|\widehat f(k)|^2 .
\label{eq:besov-difference-fourier}
\end{equation}
若$\omega_r(f,t)_2\leq\epsilon(t)t^r$且$\epsilon(t)\to0$，固定$k\neq0$，取$h=t$，只保留第$k$项，得$(2|\sin(\pi kt)|/t)^{2r}|\widehat f(k)|^2\leq\epsilon(t)^2$。令$t\to0$，左端趋于$(2\pi|k|)^{2r}|\widehat f(k)|^2$，所以$\widehat f(k)=0$。于是$f$是常数。

这个现象称为饱和：$r$阶连续模至多能表达$r$阶的规律。若想用$\omega_r(f,t)_p\leq Ct^s$刻画$s$阶光滑性，而$s>r$，那么只有常数满足条件。例如$\sin(2\pi x)$在$L^2$中有$\omega_1(\sin(2\pi\cdot),t)_2=\sqrt2\sin(\pi t)$（$t\leq1/2$），与$t$同阶，$t^{-s}\omega_1$在$s>1$时无界，尽管这个函数任意光滑。$s=r$恰好处在边界上，下一节的定义也避开它。所以刻画$s$阶光滑性时，所用差分阶数$r$必须严格大于$s$；例~\ref{ex:besov-triangle-wave}中二阶差分读出一阶差分读不出的$1+1/p$阶规律，也是这个原因。
